%% ma: L12-B08-0145
%% lop: 12
%% chuong: 2
%% bai: 8
%% dang: 12.8.49
%% loai: TN4
%% muc: VD
%% dap_an: "A"
%% nguon: "CĐ2.3.pdf giáo viên gửi 27/10/2026 (M07, Lê Minh Khôi – Nguyễn Anh Khôi, Chương II Chủ đề 3 Biểu thức tọa độ của các phép toán vectơ), Câu 61"
%% kien_thuc: "Bài 8 (điều kiện vuông góc, biểu thức tọa độ), lớp 10 (hệ phương trình bậc nhất)"
%% trang_thai: cho-duyet
%% ly_do: "Dạng mới 12.8.49 đề xuất, chờ giáo viên duyệt; Đề gốc không có phương án đúng; đã thay phương án A bằng X(3/2; 5/2; 1)"
%% ghi_chu: "Đề gốc không có phương án đúng (trực tâm đúng là (3/2; 5/2; 1)); đã thay phương án A (13/14; 27/14; 1) bằng giá trị đúng"
\begin{ex}
Trong không gian $Oxyz$, cho ba điểm $A(-1;2;3)$, $B(1;0;-1)$ và $C(2;3;1)$. Tọa độ trực tâm $X$ của tam giác $ABC$ là
\choice{\True $X\left(\dfrac32;\dfrac52;1\right)$}{$X\left(1;\dfrac{27}{14};\dfrac{13}{14}\right)$}{$X\left(-\dfrac{13}{14};\dfrac{27}{14};-1\right)$}{$X\left(\dfrac{13}{14};\dfrac{27}{14};-1\right)$}
\loigiai{$\overrightarrow{AX}=s\overrightarrow{AB}+t\overrightarrow{AC}$, $\overrightarrow{AB}=(2;-2;-4)$, $\overrightarrow{AC}=(3;1;-2)$, $\overrightarrow{BC}=(1;3;2)$.

$\overrightarrow{AX}\cdot\overrightarrow{BC}=0$, $\overrightarrow{BX}\cdot\overrightarrow{AC}=0\Rightarrow X\left(\dfrac32;\dfrac52;1\right)$ (kiểm lại: $\overrightarrow{AX}=\left(\dfrac52;\dfrac12;-2\right)\perp\overrightarrow{BC}$, $\overrightarrow{BX}=\left(\dfrac12;\dfrac52;2\right)\perp\overrightarrow{AC}$).}
\end{ex}
